                                          IMC 2022
                               Second Day, August 4, 2022
                                            Solutions


Problem 5. We colour all the sides and diagonals of a regular polygon P with 43 vertices either
red or blue in such a way that every vertex is an endpoint of 20 red segments and 22 blue segments.
A triangle formed by vertices of P is called monochromatic if all of its sides have the same colour.
Suppose that there are 2022 blue monochromatic triangles. How many red monochromatic triangles
are there?
                                                       (proposed by Mike Daas, Universiteit Leiden)

Hint: Call two connecting edges a cherry. Double-count cherries.

Solution. 1 Define a cherry to be a set of two distinct edges from K43 that have a vertex in common.
We observe that a monochromatic triangle always contains three monochromatic cherries, and that
a polychromatic triangle always contains one monochromatic cherry and two polychromatic cherries.
Therefore we study the quantity 2M − P , where M is the number of monochromatic cherries and P
is the number of polychromatic cherries. By observing that every cherry is part of a unique triangle,
we can split this quantity up into all the distinct triangles in K43 . By construction the contribution
of a polychromatic triangle will vanish, whereas a monochromatic triangle will contribute 6. We
conclude that
                        2M − P = 6 · {number of monochromatic triangles}.
Consider any vertex v. Let Mv be the number of monochromatic cherries with central vertex v and
Pv the number such polychromatic cherries. It then follows that
                             20 · 19 22 · 21
                      Mv =          +        = 421 and Pv = 20 · 22 = 440.
                                2       2
In other words, for any vertex v it holds that 2Mv − Pv = 402. Adding up all these contributions,
we find that
                                        2M − P = 43 · 402.
We conclude that there are 43 · 402/6 = 43 · 67 = 2881 monochromatic triangles in total. Since 2022
of these were blue, 859 must be red.




                                                  1
Problem 6. Let p > 2 be a prime number. Prove that there is a permutation (x1 , x2 , ..., xp−1 ) of
the numbers (1, 2, ..., p − 1) such that
                                x1 x2 + x2 x3 + ... + xp−2 xp−1 ≡ 2 (mod p).
                                       (proposed by Giorgi Arabidze, Tbilisi Free University, Georgia)

Hint:

Solution 1. We show such a permutation.
   Let xi ≡ i−1 (mod p) for i = 1, 2, · · · , p − 1. Then
          p−2           p−2         p−2          
          X             X   1  1    X     1    1         1    p−2
              xi xi+1 ≡      ·    ≡         −       ≡1−     ≡     ≡2                      (mod p)
          i=1           i=1
                            i i+1   i=1
                                          i   i+1       p−1   p−1


Solution 2. We begin by noting that the identity permutation yields the value
                                                               
                                                               p
                 1 · 2 + 2 · 3 + · · · + (p − 2)(p − 1) = 2 ·     ≡ 0 (mod p)
                                                               3

as soon as p > 3. The idea now is to perturb that permutation to obtain the desired value 2.
    One thing we can do is to replace (i, i + 1, i + 2, i + 3) by (i, i + 2, i + 1, i + 3). Indeed, this will
decrease the sum by 3. So if p ≡ 2 (mod 3), we can just take the permutation (1, 3, 2, 4, 6, 5, 7, . . . , p−
4, p−2, p−3, p−1) i.e. exchanging 3k −1 and 3k whenever k = 1, 2, . . . , p−2  3
                                                                                  . This means we decrease
           p−2
the sum 3 times by 3, leading to a remaining sum of −(p − 2) ≡ 2 (mod p).
    If p ≡ 1 (mod 3), this strategy does not work immediately. Instead, we can change (1, 2, 3, 4, 5)
to (1, 4, 3, 2, 5) resulting in a decrement of the sum by 8. If we then exchange 3k and 3k + 1 for
k = 2, 3, . . . , p−7
                    3
                      as before, we get another p−10
                                                 3
                                                     times a decrement by 3, leading to a remaining sum
of −8 − p−10 3
                  · 3 ≡ 2 (mod  p).
    Of course this only works if p ≥ 13. It thus remains to consider the cases p = 3 and p = 7 by
hand. For p = 3, we just take (1, 2) and for p = 7 we can take (1, 4, 5, 2, 3, 6).




                                                     2
Problem 7. Let A1 , A2 , . . . , Ak be n × n idempotent complex matrices such that

                                           Ai Aj = −Aj Ai                for all i ̸= j.

Prove that at least one of the given matrices has rank ≤ nk .
(A matrix A is called idempotent if A2 = A.)
                                                                              (proposed by Danila Belousov, Novosibirsk)

Hint: Consider the trace and the rank of A.

Solution 1.
Lemma. For any idempotent matrix B

                                                         tr(B) = rank(B)

Proof. Observe that an idempotent matrix satisfies the equation λ(1 − λ) = 0. Hence the minimal
polynomial is a product of linear factors and the matrix is diagonalizable. Therefore, the rank of
the matrix equals the number of non-zero eigenvalues. Since the matrix has eigenvalues 0 or 1, this
provides that the trace is equal to the number of unity eigenvalues, or non-zero eigenvalues.
                          Pk
   It can be shown that      Ai is also an idempotent. Indeed,
                             i=1

                              k
                             X           2       k
                                                   X             X                              k
                                                                                                X
                                     Ai        =         A2i +          (Ai Aj + Aj Ai ) =            Ai
                               i=1                 i=1           i̸=j                           i=1


   Applying the lemma one can obtain
                    k
                    X                      k
                                           X                        k
                                                                   X                           k
                                                                                                X         
                           rank(Ai ) =             tr(Ai ) = tr               Ai       = rank         Ai       ⩽n
                     i=1                   i=1                          i=1                     i=1

   The required inequality follows.

Solution 2. We first prove that for idempotents A, B with AB = −BA we already must have
AB = BA = 0. Indeed, it is clear that ABx = BAx = 0 for x ∈ ker(A) so it suffices to prove
the same for x ∈ im(A), i.e. when Ax = x. But then writing Bx = y we have Ay = −y i.e.
y = −Ay = −A2 y = Ay = −y and hence y = 0 so that again ABx = BAx = 0.
    Henceforth, we can assume the stronger condition Ai Aj = 0 for all i ̸= j. We next claim that
all the image spaces Vi of Ai are linearly independent. This will imply the claim, since then the
sum of their dimensions canPbe at most n, and so one of them has to be ≤ nk . Now, for the sake of
contradiction, suppose that i vi = 0 with vi ∈ Vi and w.l.o.g. v1 ̸= 0. But then

         0 = A1 (v1 + · · · + vk ) = v1 + A1 v2 + · · · + A1 vk = v1 + A1 A2 v2 + · · · + A1 Ak vk = v1

since A1 Ai = 0 for all i.

Remark. Here is a different argument for AB = BA = 0, without eigenvectors: multiplying by A and using
its idempotence and the super-commutativity , we have

                        −BA = AB = A2 B = AAB = −ABA = BAA = BA2 = BA

thus BA = 0.


                                                                  3
Problem 8. Let n, k ≥ 3 be integers, and let S be a circle. Let n blue points and k red points be
chosen uniformly and independently at random on the circle S. Denote by F the intersection of the
convex hull of the red points and the convex hull of the blue points. Let m be the number of vertices
of the convex polygon F (in particular, m = 0 when F is empty). Find the expected value of m.
                                                                 (proposed by Fedor Petrov, St. Petersburg)

Hint:

Solution 1. We prove that
                                                 2kn         k!n!
                                     E(m) =           −2              .
                                                n+k−1    (k + n − 1)!

   Let A1 , . . . , An be blue points. Fix i ∈ {1, . . . , n}. Enumerate our n + k points starting from a
blue point Ai counterclockwise as Ai , X1,i , X2,i , . . . , X(n+k−1),i . Denote the minimal index j for which
the point Xj,i is blue as m(i). So, Ai Xm(i),i is a side of the convex hull of blue points. Denote by bi
the following random variable:
                                 (
                                  1, if the chord Ai Xm(i),i contains a side of F
                            bi =
                                  0, otherwise.

Define analogously k random variables r1 , . . . , rk for the red points. Clearly,

                                  m = b1 + . . . + bn + r1 + . . . + rk .     (♡)

    We proceed with computing the expectation of each bi and rj . Note that bi = 0 if and only if all
red points lie on the side of the line Ai Xm(i),i . This happens either if m(i) = 1, i.e., the point Xi,1 is
                                            n−1
blue (which happens with probability k+n−1      ), or if i = k + 1, points X1,i , . . . , Xk,i are red, and points
Xk+1,i , . . . , Xk+n−1,i are blue (which happens with probability 1/ k+n−1
                                                                                 
                                                                            k
                                                                                   , since all subsets of size k
of {1, 2, . . . , n + k − 1} have equal probabilities to correspond to the indices of red points between
                                                                                                           k!(n−1)!
                                                                     n−1
X1,i , . . . , Xn+k−1,i ). Thus the expectation of bi equals 1 − k+n−1    − 1/ k+n−1  k
                                                                                                   k
                                                                                             = n+k−1    − (k+n−1)!  .
                                            n      n!(k−1)!
Analogously, the expectation of rj equals n+k−1 − (k+n−1)!  . It remains to use (♡) and linearity of
expectation.

Solution 2. Let C1 , . . . , Cn+k be the colours of the points, scanned counterclockwise from a fixed
point on the circle. We consider the sequence as cyclic (so Cn+k is also adjacent to C1 ). There are
two cases: Either (i) all red points appear contiguously, followed by all blue points contiguously, or
(ii) the red and blue points alternate at least twice. It can be seen that in the second case, m is
exactly equal to the number of colour changes in the Ci sequence: For example, if Ci is red and Ci+1
is blue, then the intersection of the red chord from Ci to the next red point with the blue chord
from Ci+1 to the previous blue point is a vertex of F , and every vertex is of this form. Case (i) is
exceptional, as we have two colour changes, but m = 0, so it is 2 less than the number of changes in
that case.
    Now observe that the distribution of Ci is purely combinatorial: Each of the n+k
                                                                                            
                                                                                        n,k
                                                                                              distributions
of colours is equally likely (for example, because we can generate the distribution by first choosing all
n + k points on the circle, and then assigning colours uniformly). In particular the probability that
                                             2nk
Ci Ci+1 is a colour change is exactly (n+k)(n+k−1)  , and by lineraity of expectation, the total expected
                                                                             2nk
number of color changes (including i = n + k) is n + k times this, i.e. n+k−1     .
    To get the expected value of m, we must subtract from the above 2 times the probability of case
                                                                                                      −1
(i). Exactly n + k of the n+k     distributions belong to case (i), so we must subtract 2(n + k) n+k
                                
                            n,k                                                                   n,k
                                                                                                          =
    n!k!
2 (n+k−1)! , as claimed.

                                                         4
Solution 3. Let A1 , . . . , An be the blue points and B1 , . . . , Bk be the red points. For every pair of
blue points Ai , Aj , 1 ⩽ i < j ⩽ n, we evaluate the probability p that Ai Aj contains a side of F (it
obviously does not depend on the choice of i and j). By q denote      the    analogous probability for the
                                                                    n      k
red points. Then by linearity of expectation we have Em = 2 p + 2 q.
    We proceed with finding p. Without loss of generality i = 1, j = 2. Let the length of the circle be
1, and the length of arc A1 A2 (counterclockwise from A1 to A2 ) be x. Then x is uniformly distributed
on [0, 1]. Then A1 A2 contains a side of F if
    (i) all blue points are on the same side of A1 A2 , but
    (ii) the red points are not on the same side of A1 A2 .
    The probability of (i) is xn−2 + (1 − x)n−2 . The probability of (ii) is 1 − (xk + (1 − x)n−k ).
                                                 R1
    Thus, using Beta function value B(a, b) = 0 xa−1 (1 − x)b−1 dx = B(a, b) = (a−1)!(b−1)!
                                                                                       (a+b−1)!
                                                                                                for positive
integers a, b
         Z 1
                                                                    2           2
    p=       (xn−2 + (1 − x)n−2 )(1 − (xk + (1 − x)n−k ))dx =           −             − 2B(n − 1, k + 1)
          0                                                     n−1 n+k−1
            2          2          (n − 2)!k!
      =         −            −2               .
         n−1 n+k−1               (n + k − 1)!

Next,             
                   n        n(n − 1)        n!k!            nk          n!k!
                      p=n−           −                =           −             ,
                   2        n + k − 1 (n + k − 1)!       n + k − 1 (n + k − 1)!
and by symmetry k2 q takes the same value (that is in agreement with the observation that red and
                     

blue sides of F alternate).




                                                     5
